\documentclass[preview, border={10mm, 5mm, 10mm, 10mm}]{standalone}
\usepackage{tikz, tasks, geometry, xcolor}
\usetikzlibrary{arrows.meta, patterns, calc, intersections, quotes, angles}
\usepackage{amsmath, amssymb, amsfonts, mathtools}
\setlength{\columnsep}{10pt}
\setlength{\columnseprule}{0.4pt}
\usepackage[upright]{fourier}
\usepackage{enumitem}
\geometry{a4paper, margin=1in}
\usepackage{tzplot, pgfplots, kinematikz}
\usepackage{circuitikz}
\ctikzset{resistors/scale=0.75,capacitors/scale=0.75,inductors/scale=0.75}
\usepackage{chemfig}
\usepackage[version=4]{mhchem}
\usepackage[modules=all]{chemmacros}
\usepackage{tkz-euclide}
\usepackage{venndiagram}
\pgfplotsset{compat=1.18}
\everymath{\displaystyle}
\newcommand{\ans}{\textcolor{blue!20!red}{\textit{\quad Ans.}}}
% \renewcommand{\ans}{}  % Uncomment to hide answers
\newenvironment{solution}{\par\noindent\color{red!80!black}$\Rightarrow$\enspace\ignorespaces}{\par}
\newenvironment{alternatesolution}{\par\noindent\color{blue!80!black}$\Rrightarrow$\enspace\ignorespaces}{\par}
\newenvironment{hint}{\par\noindent\color{red!50!black}$\looparrowright$\enspace\ignorespaces}{\par}
\newenvironment{idea}{\par\noindent\color{violet!80!black}$\diamond$\enspace\ignorespaces}{\par}
\newenvironment{remark}{\par\noindent\color{teal!80!black}$\circ$\enspace\ignorespaces}{\par}

% --- Global TikZ style (design uniformity across all diagrams) ---
\tikzset{
    >=latex,
    thick,
    every node/.append style={font=\small},
}
\title{\textsc{}}
\pagestyle{empty}
\begin{document}
\begin{solution}
\begin{align*}
v &= \frac{c}{100} \\
  &= \frac{3 \times 10^8}{100} \\
  &= 3 \times 10^6 \ \mathrm{m s}^{-1} \\
\intertext{For no deflection, net Lorentz force must be zero. Hence}
q\vec{E} + q\left( \vec{v} \times \vec{B} \right) &= 0 \\
\vec{E} &= -\left( \vec{v} \times \vec{B} \right) \\
\intertext{Thus $\vec{E}$ must be along the direction opposite to $\vec{v} \times \vec{B}$, so it is perpendicular to $\vec{B}$. Its magnitude is}
E &= v B \\
  &= \left( 3 \times 10^6 \right)\left( 9 \times 10^{-4} \right) \\
  &= 27 \times 10^2 \ \mathrm{V m}^{-1} \\
\intertext{So $\vec{E}$ is perpendicular to $\vec{B}$ and has magnitude $27 \times 10^2\,\mathrm{V m}^{-1}$. Therefore, the correct option is (d).}
\end{align*}
\end{solution}
\hfill $\_\_$
\end{document}
